\documentclass[12pt, openany]{book}

% ──────────────────────────────────────────────
% Packages
% ──────────────────────────────────────────────
\usepackage[margin=1in]{geometry}
\usepackage{amsmath, amssymb, amsthm}
\usepackage{mathtools}
\usepackage{bm}
\usepackage{tikz}
\usepackage{pgfplots}
\pgfplotsset{compat=1.18}
\usetikzlibrary{arrows.meta, decorations.markings, calc, patterns, positioning}
\usepackage{tcolorbox}
\tcbuselibrary{theorems, breakable, skins}
\usepackage{hyperref}
\usepackage{enumitem}
\usepackage{fancyhdr}
\usepackage{titlesec}
\usepackage{epigraph}
\usepackage{xcolor}
\usepackage{booktabs}

% ──────────────────────────────────────────────
% Colors
% ──────────────────────────────────────────────
\definecolor{chapblue}{RGB}{0, 70, 130}
\definecolor{accentred}{RGB}{180, 40, 40}
\definecolor{softgray}{RGB}{245, 245, 245}
\definecolor{trajgreen}{RGB}{30, 130, 60}
\definecolor{darkgold}{RGB}{160, 120, 20}
\definecolor{purplehaze}{RGB}{100, 40, 140}

% ──────────────────────────────────────────────
% Theorem environments
% ──────────────────────────────────────────────
\newtcbtheorem[number within=chapter]{principle}{Principle}{
  colback=chapblue!5, colframe=chapblue!80,
  fonttitle=\bfseries, separator sign={.~},
  breakable
}{princ}

\newtcbtheorem[number within=chapter]{keyidea}{Key Idea}{
  colback=accentred!5, colframe=accentred!70,
  fonttitle=\bfseries, separator sign={.~},
  breakable
}{idea}

\newtcbtheorem[number within=chapter]{example}{Example}{
  colback=softgray, colframe=black!40,
  fonttitle=\bfseries, separator sign={.~},
  breakable
}{ex}

\newtcbtheorem[number within=chapter]{definition}{Definition}{
  colback=darkgold!5, colframe=darkgold!80,
  fonttitle=\bfseries, separator sign={.~},
  breakable
}{def}

\newtcbtheorem[number within=chapter]{remark}{Remark}{
  colback=purplehaze!5, colframe=purplehaze!60,
  fonttitle=\bfseries, separator sign={.~},
  breakable
}{rem}

\newtcbtheorem[number within=chapter]{warning}{Warning}{
  colback=accentred!5, colframe=accentred!80,
  fonttitle=\bfseries, separator sign={.~},
  breakable
}{warn}

\theoremstyle{plain}
\newtheorem{proposition}{Proposition}[chapter]
\newtheorem{theorem}{Theorem}[chapter]
\newtheorem{lemma}[theorem]{Lemma}
\newtheorem{corollary}[theorem]{Corollary}

% ──────────────────────────────────────────────
% Chapter formatting
% ──────────────────────────────────────────────
\titleformat{\chapter}[display]
  {\normalfont\huge\bfseries\color{chapblue}}
  {\chaptertitlename\ \thechapter}{20pt}{\Huge}
\titlespacing*{\chapter}{0pt}{-20pt}{30pt}

% ──────────────────────────────────────────────
% Header/Footer
% ──────────────────────────────────────────────
\pagestyle{fancy}
\fancyhf{}
\fancyhead[LE,RO]{\thepage}
\fancyhead[RE]{\textit{Control Is Motion}}
\fancyhead[LO]{\textit{\nouppercase{\leftmark}}}
\renewcommand{\headrulewidth}{0.4pt}
\setlength{\headheight}{14.5pt}

% ──────────────────────────────────────────────
% Custom commands
% ──────────────────────────────────────────────
\newcommand{\state}{\bm{x}}
\newcommand{\control}{\bm{u}}
\newcommand{\traj}{\bm{\gamma}}
\newcommand{\manifold}{\mathcal{M}}
\newcommand{\configspace}{\mathcal{Q}}
\newcommand{\statespace}{\mathcal{X}}
\newcommand{\controlspace}{\mathcal{U}}
\newcommand{\Reals}{\mathbb{R}}
\newcommand{\dd}{\mathrm{d}}
\newcommand{\dt}{\,\dd t}
\newcommand{\norm}[1]{\left\lVert #1 \right\rVert}
\newcommand{\inner}[2]{\left\langle #1,\, #2 \right\rangle}

% ──────────────────────────────────────────────
\begin{document}

\setcounter{chapter}{8}

\chapter{Stochastic Trajectories \& Motor Variability}

\epigraph{%
  Perfection is inherently fragile. Biological optimization seeks the strongest balance between intent and reality.
}{}

\vspace{0.5cm}

% ══════════════════════════════════════════════
\section{Signal-Dependent Noise}
% ══════════════════════════════════════════════

When a roboticist specifies a torque of $50\text{ Nm}$, the motor delivers $50\text{ Nm}$ with a tiny, constant variance. Biological actuators do not work this way. Human muscles generate force by recruiting discrete motor units. As higher forces are required, larger, more fast-twitch oriented motor units are recruited. This physiological reality manifests as a profound control challenge: \textbf{motor noise is proportional to the command signal}.

\begin{definition}{Signal-Dependent Noise (SDN)}{sdn}
  A control input $\control(t)$ applied by a biological system is actually realized as a stochastic variable $\tilde{\control}(t)$:
  \begin{equation}
    \tilde{\control}(t) = \control(t) + \bm{W} \control(t) \cdot \bm{\epsilon}(t)
  \end{equation}
  where $\bm{W}$ is a scaling matrix representing the coefficient of variation, and $\bm{\epsilon}(t) \sim \mathcal{N}(0, I)$ is white noise. That is, the variance of the applied force grows squarely with the magnitude of the ordered force: $\text{Var}(\tilde{u}) = c u^2$.
\end{definition}

This simple reality crumbles deterministic optimal control. If you compute an optimal trajectory that requires maximum muscular exertion, the trajectory will mathematically possess minimum variance only if noise is constant. Under SDN, maximum exertion guarantees maximum chaotic noise.

% ══════════════════════════════════════════════
\section{Minimum Variance Theory of Human Movement}
% ══════════════════════════════════════════════

Why do humans move in smooth, bell-shaped velocity profiles rather than jerky "bang-bang" optimal control solutions when reaching for coffee?

Classical control attempts to explain this via "minimum jerk" theories, positing that the nervous system explicitly encodes an optimization functional to minimize $\int \dddot{x}^2 \dt$. The trajectory-first perspective offers a much deeper, physically grounded explanation championed by Harris and Wolpert: the \textbf{Minimum Variance Theory}.

\begin{keyidea}{Optimization for Repeatability}{min_variance}
  The human nervous system seeks to move in a way that minimizes the final endpoint variance evaluated across multiple trials.
\end{keyidea}

Under SDN, the only way to minimize endpoint variance (ensure you actually grasp the coffee cup smoothly) is to minimize sudden spikes in required actuator forces, because large forces inject huge amounts of uncertainty into the dynamics. Smooth trajectories inherently possess lower peak torque commands than rapid accelerations and decelerations. Smoothness is not the explicitly programmed objective; it is an emergent property of attempting to be consistently accurate under the unavoidable reality of signal-dependent noise.

% ══════════════════════════════════════════════
\section{The Speed-Accuracy Tradeoff in the Golf Swing}
% ══════════════════════════════════════════════

Fitts's Law states that the faster you try to move to a target, the less accurate you become. In golf, we see the supreme manifestation of this. The deterministic optimal control solution (Chapter 6) tells the solver to maximize input torque to achieve theoretically maximal clubhead speed. 

However, if we introduce the Stochastic OCP (via Covariance Steering) from Chapter 6 and apply the realistic premise of SDN, the optimizer faces a profound tradeoff.

Applying maximum $u_1$ torque during the downswing injects enormous variance into the kinematic state. The unactuated club ($\theta_2$) begins whipping through its zero dynamics, but its exact phase progression is now highly uncertain due to the noise flooding the arms. 

To ensure the clubface squares up at the moving goal, the stochastic optimizer will \emph{deliberately back off} the peak requested torque. It trades a deterministic $5$ mph gain in clubhead speed for a massive reduction in angular variance at the target manifold. The optimal biological swing is mathematically proven to be a sub-maximal physical exertion.

\end{document}
